\chapter{Definitely not the same thing as optimism (I Won't Rely on Anyone Anymore)}
Let's start back at the beginning.

Among many other things, Madoka Magica is a critique of moe, the fandom within anime fandom that celebrates weak, usually feminine characters that evoke a feeling of protectiveness. As we saw in Sayaka's arc, this protectiveness is futile and dangerous, because to be protected is to be objectified, to be treated as other than a full human being possessed of free agency.

Just as magical girls are witches, to be a protector is to be a destroyer, and this is just as true of Homura (whose family name, Akemi, means "fire," suggesting a connection to the "burning" of moe) as Sayaka. Note the precise nature of Homura's wish: to switch places with Madoka, to be the one to protect her instead of protected by her. By setting herself up as Madoka's protector, she inflicts the status of protected on Madoka, which (as we learn next episode) is the very reason Madoka has suffered so much—but that should not be a surprise, as over-protecting someone can very easily lead to them losing confidence.

The series' Buddhist roots are particularly important here, because Homura is being trapped by her own conception—rather than consult what Madoka wants, she is imposing her own notions of what she believes Madoka should want, which is really based in Homura's own desires. Homura's desires thus lead her to suffering, as the Buddha taught that they would \cite{ref75}. She is also being led astray by believing that the concepts she imposes on the world are themselves real; by contrast, the impermanent and ever-changing world (including the impacts of her own wishes and actions) denies the distinctions she tries to draw, such as future/past, alive/dead, and human/magical girl/witch. This is very much the Zen notion that language and thought interfere with our ability to see the panenhenic truth that, rather than consisting of independent entities and states, the universe is only the ever-fluxing one-in-all \cite{ref76}. 

Homura first struggles to save Madoka from death because she sees it as different from life; then she struggles to save Madoka from becoming a witch because she sees it as different from being a magical girl; finally she struggles to save Madoka from becoming a magical girl because she sees it as different from being an ordinary human. But because she fundamentally does not understand what kind of a story she's in—and it is telling that she went to a Christian school before coming to Mitakihara Middle School—she cannot recognize that her actions will only ever make things worse.

In short, this isn't working.

Let's start back at the beginning.

Last chapter I discussed the magical girl genre's origins, and how the good girl/bad girl (or magical girl/witch, or Madonna/whore) dichotomy constrains the expression of feminine power into either channels non-threatening to masculine hegemony or into outsider figures, whether the witch deep in the woods or the magical girl who only ever fights monsters no one else can see.

Episode 10 uses multiple timelines to spread Homura's character out, and show her evolution from the non-threatening form to the outsider. In the first timeline, she is a powerless innocent, weak enough to become a witch's victim—although, foreshadowing her eventual cold determination, despite her physical frailty she is the only witch victim in the series who does not actually succumb to the witch's mental influence. As the timelines progress, she becomes more powerful and more of a transgressive figure. At the beginning of the second timeline, she is clearly the weakest and least skilled of the magical girls; her power to stop time is purely a support ability with no apparent offensive application, and her ability to physically harm witches is minimal. By the end of the timeline, however, she is capable of making and using homemade bombs in combination with her time stop to kill witches, and Madoka is delegated to the support role. In the third timeline, Homura becomes more powerful and more transgressive still, moving from making bombs at home to stealing guns from criminals, making her a match for Oktavia.

In the fourth timeline, Homura finally becomes the figure we know, coldly detached, her girlish pigtails and glasses—both signifiers of harmlessness—abandoned along with her uncertainty, as she robs the ultimate symbol of hegemonic masculinity, the military, to fight alone.

But these images of Homura are spread out, not from past to future, but across multiple timelines, meaning that in some sense all these Homuras exist simultaneously. The "outside" in which the outsider fights is a channel non-threatening to masculine hegemony. The apparent binary between them is just as illusory as that between magical girl and witch; Mami and Homura are two sides of the same coin.

Which means that Homura is not a rebuke of the magical girl genre or a challenge to the safe, non-challenging framework within which it resides. She is simply another manifestation of that framework, an alternate form of the magical girl with whom viewers are comfortable.

In short, this isn't working.

Let's start back at the beginning.

In Chapter 1, I posited that there are three arcs to Madoka Magica, each focused on a particular character or pair of characters and each emphasizing different themes. Episode 9 closed out the Sayaka/Kyoko arc, which was largely about exploring depression and closing the circle on the magical girl genre. Episode 10, thus, serves as an introduction to the final arc, which focuses primarily on Homura.

We thus get an episode dedicated entirely to her perspective, following her as she repeatedly travels back in time in an attempt to avert Madoka's fate. Throughout this, Homura appears largely immune to the despair that overwhelmed Sayaka in her arc; the only time Homura comes close to becoming a witch in this episode is in the third timeline, when she suggests that she and Madoka become monsters and wipe away the entire world and all its sadness, which foreshadows both Madoka's eventual wish and much of the plot of Rebellion.

What makes the third timeline special? In the first timeline, Homura only becomes a magical girl after Walpurgisnacht kills Madoka. In the second, Homura's lack of injury suggests that she had little or no part in the fight with Walpurgisnacht. In the fourth, Madoka only becomes a magical girl after Walpurgisnacht defeats Homura.

The timeline on which Homura comes closest to being a witch is the timeline on which she actually gets her wish, to defend Madoka during the battle with Walpurgisnacht. This is the balance of hope and despair which Sayaka spoke about; the heights lead directly to the depths.

Homura is usually immune to despair, because she does not hope in the usual sense. Instead of hope, she relies on her determination and her love for Madoka; since neither is an innately good feeling, neither brings despair to balance it. Put another way, she epitomizes the concept of hope as described by the Czech playwright, philosopher, revolutionary, and statesman Vaclav Havel: "[hope] is not the conviction that something will turn out well, but the certainty that something makes sense, regardless of how it turns out." \cite{ref77} 

As she says at the end of the episode, she can keep going no matter how long it takes, so long as it is for Madoka. But what then to make of her breakdown in Episode 8? Madoka seems determined to sacrifice herself in every timeline, and this is the single greatest threat to Homura's plan. Homura is functionally depressed, which is to say that while she is functional, she is also depressed, and all it takes is a slight push downward to make her collapse.

In short, this isn't working.

Let's start back at the beginning.

Episode 10 presents us with a series of alternate timelines, adjacent to the events of the show. As each timeline reaches the apocalyptic disaster of Walpurgisnacht, Homura resets herself to the day she left the hospital and tries again. However, it appears that she is not just traveling back in time, but also across timelines, as elements seem to change from timeline to timeline regardless of Homura's involvement—for example, in the first timeline Madoka becomes a magical girl before Homura leaves the hospital.

Throughout these four alternate timelines, there are echoes (or omens, depending on your point of view) of events in the series' main timeline, from the shot of Madoka's shoes as she turns to face Homura in the school hallway in the first timeline (an echo of a similar shot of Homura turning in the same scene in the first episode) to Kyoko's refusal to accept that Oktavia has completely replaced the Sayaka she knew.

But more interesting are the subtler echoes. Homura's repeated awakenings in the hospital, for instance, echo Charlotte's Witch Seed hatching at a hospital in Episode 3, strengthening the parallel between them as harbingers of the true series and Mami's nemeses. Madoka's pleading with Homura to prevent her from making the mistake of contracting with Kyubey places Homura in the position of making that mistake in Madoka's place, echoing Junko's advice in Episode 6. Even Homura's gradual loss of innocence and transformation into a darker, more mature, and more transgressive figure parallels Sayaka's entire arc.

These echoes serve not only to make clear how the timelines relate to one another; they also recontextualize the events they echo. From Homura's point of view, these are not echoes but prefigurations (as, indeed, some are for the viewer as well, such as Mami being the first to attack the other magical girls, foreshadowing Rebellion); all this has happened before, and likely will happen again.

Two episodes utterly transform the series. On a first viewing, these two episodes stand out as moments when what appeared to be a show about one thing suddenly becomes a different show entirely. As in a mirror, they are at precisely opposite ends of the series, the third episode and the third-to-last. This is the second of these episodes.

If nothing else (and it is much else), Puella Magi Madoka Magica is a meticulously structured series. Every moment of it is carefully placed to advance a complex story with extensive character development in a surprisingly small space. One example of that meticulous structure is the symmetry throughout the series, the way in which moments large and small repeat at precisely the right time.

In short, this isn't working.
